\chapter{面积与体积}\label{ch:g7:areas}

六年级学习了长方形、直角三角形和长方体的度量（见\cref{ch:g6:measure}）。
用一把剪刀——剪下这里的一块，粘到那里——
同样的想法，现在可以求出平行四边形、
\emph{所有}三角形和圆的面积，以及棱柱和
圆柱的体积。

\section{平行四边形的面积}

\begin{theorem}[平行四边形]\label{thm:g7:areas:parallelogram}
设平行四边形的一边长为 $b$（称其为\emph{底}），
该边与其对边之间的距离为 $h$（称其为
\emph{高}），则其面积为
\[
A = b \times h .
\]
\end{theorem}

\begin{proof}[剪拼证明]
把伸出一侧的直角三角形剪下，
粘到另一侧，平行四边形就变成了一个
面积不变的 $b \times h$ 长方形。
\end{proof}

\begin{omfigure}
\begin{tikzpicture}[scale=0.82]
  \fill[omDef!14] (0,0) -- (4,0) -- (5.2,1.8) -- (1.2,1.8) -- cycle;
  \draw[thick] (0,0) -- (4,0) -- (5.2,1.8) -- (1.2,1.8) -- cycle;
  \fill[omThm!25] (0,0) -- (1.2,1.8) -- (1.2,0) -- cycle;
  \draw[omThm, thick] (0,0) -- (1.2,1.8) -- (1.2,0) -- cycle;
  \draw[dashed] (1.2,0) -- (1.2,1.8);
  \node[below, font=\small] at (2.0,0) {$b$};
  \draw[Stealth-Stealth] (6.0,0) -- (6.0,1.8)
    node[midway, right, font=\small] {$h$};
  \begin{scope}[xshift=8.2cm]
  \fill[omDef!14] (0,0) rectangle (4,1.8);
  \draw[thick] (0,0) rectangle (4,1.8);
  \fill[omThm!25] (2.8,0) -- (4,1.8) -- (4,0) -- cycle;
  \draw[omThm, thick] (2.8,0) -- (4,1.8) -- (4,0) -- cycle;
  \node[below, font=\small] at (2.0,0) {$b$};
  \end{scope}
\end{tikzpicture}

{\small 剪拼证明：把红色三角形从左端平移到右端，
平行四边形就变成了长方形——面积相同，
都是 $b \times h$。}
\end{omfigure}

\begin{remark}
高是两条底之间的\emph{垂直}距离，
而不是斜边的长度！非常倾斜的平行四边形，
边可以很长，而面积却很小。
\end{remark}

\section{三角形的面积}

\begin{theorem}[三角形]\label{thm:g7:areas:triangle}
设三角形的底为 $b$，对应的高为 $h$（即
对顶点到底边所在直线的垂直距离），
则其面积为
\[
A = \frac{b \times h}{2} .
\]
\end{theorem}

\begin{proof}
两个相同的三角形，其中一个绕某一边的中点旋转半圈，
与另一个拼成一个底为 $b$、
高为 $h$ 的平行四边形（这正是中心对称的作用，
见\cref{prop:g7:central:preserved}）。三角形恰为其一半，
即 $\frac{bh}{2}$。
\end{proof}

\begin{omfigure}
\begin{tikzpicture}[scale=0.9]
  \fill[omDef!14] (0,0) -- (4,0) -- (1.4,2.1) -- cycle;
  \draw[thick] (0,0) -- (4,0) -- (1.4,2.1) -- cycle;
  \fill[omThm!12] (4,0) -- (1.4,2.1) -- (5.4,2.1) -- cycle;
  \draw[thick, omThm, dashed] (4,0) -- (5.4,2.1) -- (1.4,2.1);
  \draw[dashed] (1.4,0) -- (1.4,2.1);
  \draw[thin] (1.4,0) ++(-0.2,0) -- ++(0,0.2) -- ++(0.2,0);
  \node[below, font=\small] at (2,0) {$b$};
  \node[right, font=\small] at (1.42,1.0) {$h$};
\end{tikzpicture}

{\small 三角形与其旋转半圈所得的副本（虚线）拼成一个
平行四边形：三角形的面积为 $\frac{b\times h}{2}$。三条边中的任意一条都可以作底，
但必须配\emph{自己}对应的高。}
\end{omfigure}

\begin{example}\label{ex:g7:areas:triangle}
一个三角形的底为 $9$ cm，
对应的高为 $4$ cm：
\[
A = \frac{9 \times 4}{2} = \frac{36}{2} = 18 \text{ cm}^2 .
\]
对于\emph{直角}三角形，两条直角边恰好是一组底和高：
由此可得到\cref{prop:g6:measure:formulas}中的公式。
\end{example}

\section{圆的面积}

\begin{theorem}[圆]\label{thm:g7:areas:disk}
半径为 $r$ 的圆，其面积为
\[
A = \pi r^2 .
\]
\end{theorem}

\begin{proof}[证明思路]
将圆盘切成许多等份的细扇形，首尾相接排列：
它们近似拼成一个高为 $r$、底为半个圆周的
平行四边形，底长为 $\pi r$（见\cref{prop:g6:measure:circle}）。其面积
近似为 $\pi r \times r = \pi r^2$，且切片越细，近似越精确。
完整的证明需要用到积分学，
这部分内容见高中分册。
\end{proof}

\begin{omfigure}
\begin{tikzpicture}[scale=0.85]
  \foreach \a in {0,45,...,315}
    \fill[omDef!20, draw=omDef] (0,0) -- (\a:1.4)
      arc (\a:\a+45:1.4) -- cycle;
  \begin{scope}[xshift=4.3cm, yshift=-0.15cm]
  \foreach \i in {0,1,2,3} {
    \fill[omDef!20, draw=omDef]
      ({\i*1.1},0) -- ({\i*1.1+0.55},1.4) -- ({\i*1.1+1.1},0) -- cycle;
    \fill[omDef!20, draw=omDef]
      ({\i*1.1+0.55},1.4) -- ({\i*1.1+1.1},0) -- ({\i*1.1+1.65},1.4)
      -- cycle;
  }
  \draw[Stealth-Stealth] (5.2,0) -- (5.2,1.4)
    node[midway, right, font=\small] {$\approx r$};
  \node[below, font=\small] at (2.5,-0.1) {底 $\approx \pi r$};
  \end{scope}
\end{tikzpicture}

{\small 将圆盘切开并展开各切片：近似为一个
底为 $\pi r$（半个圆周）、高为 $r$ 的平行四边形。}
\end{omfigure}

\begin{example}\label{ex:g7:areas:disk}
一张圆形桌面的半径为 $60$ cm。其面积为
$\pi \times 60^2 = 3600\pi \approx 11\,310$ cm$^2$，略大于
$1.1$ m$^2$。注意区分：\emph{周长}$2\pi r$ 只用到
$r$，而\emph{面积}$\pi r^2$ 用到的是 $r^2$。
\end{example}

\section{棱柱与圆柱}

\begin{definition}[棱柱、圆柱]\label{def:g7:areas:prism}
\emph{棱柱}\index{棱柱}是一种立体图形，它有两个全等且平行的多边形面
（\emph{底面}），这两个面之间由长方形连接；底面换成圆面时，称为
\emph{圆柱}\index{圆柱}。两底面之间的距离称为
\emph{高}。
\end{definition}

\begin{theorem}[体积]\label{thm:g7:areas:volume}
对底面积为 $B$、高为 $h$ 的棱柱或圆柱，有
\[
V = B \times h .
\]
\end{theorem}
\admitted

\begin{example}\label{ex:g7:areas:volume}
帐篷是一个侧放的棱柱：其底面为三角形，
底为 $2$ m、高为 $1.5$ m，帐篷长 $3$ m。
\begin{enumerate}
  \item 底面积：$B = \frac{2 \times 1.5}{2} = 1.5$ m$^2$。
  \item 体积：$V = B \times h = 1.5 \times 3 = 4.5$ m$^3$。
\end{enumerate}
一个半径为 $4$ cm、高为 $11$ cm 的圆柱形铁罐：
$V = \pi \times 4^2 \times 11 = 176\pi \approx 553$ cm$^3$，约
为半升。
\end{example}

\section{练习}

\begin{exercise}[$\star$]\label{exo:g7:areas:1}
求底为 $8$ cm、高为
$5$ cm 的平行四边形的面积，以及底为 $6.5$ m、高为 $4$ m 的另一个平行四边形的面积。
\end{exercise}

\begin{exercise}[$\star$]\label{exo:g7:areas:2}
一个平行四边形的两边长分别为 $10$ cm 和 $6$ cm，
相对于 $10$ cm 底边的高为 $4$ cm。求其面积。为什么
答案不是 $60$ cm$^2$？
\end{exercise}

\begin{exercise}[$\star$]\label{exo:g7:areas:3}
求底为 $12$ cm、高为 $7$ cm 的三角形面积；
直角边分别为 $9$ 和 $6$ 的直角三角形面积；
底为 $5.5$ m、高为 $2$ m 的三角形面积。
\end{exercise}

\begin{exercise}[$\star$]\label{exo:g7:areas:4}
求半径为 $5$ cm 的圆的面积，再求半径为 $10$ cm 的半圆的面积
（先给出含 $\pi$ 的精确值，再四舍五入到个位）。
\end{exercise}

\begin{exercise}[$\star$]\label{exo:g7:areas:5}
求半径为 $3$ cm 的圆的周长\emph{和}面积。
两个结果中，哪一个单位是 cm，哪一个单位是 cm$^2$？
\end{exercise}

\begin{exercise}[$\star$]\label{exo:g7:areas:6}
求底面积为 $24$ cm$^2$、高为
$10$ cm 的棱柱的体积；半径为 $3$ cm、高为 $8$ cm 的圆柱的体积（先给出
精确值，再四舍五入）。
\end{exercise}

\begin{exercise}[$\star$]\label{exo:g7:areas:7}
一个三角形的面积为 $28$ cm$^2$，底为 $8$ cm。
求对应的高，先列出方程。
\end{exercise}

\begin{exercise}[$\star\star$]\label{exo:g7:areas:8}
任意画一个三角形，仔细测量三组底与高的数值。
对每组数据计算 $\frac{b \times h}{2}$，三个结果应
在测量误差范围内相等。为什么？
\end{exercise}

\begin{exercise}[$\star\star$]\label{exo:g7:areas:9}
一块梯形花园可用一条对角线
分成两个三角形，它们的高相同，都是 $h = 20$ m，底分别为 $30$ m 和
$18$ m。求该花园的总面积。你能猜出
梯形的一般面积公式吗？
\end{exercise}

\begin{exercise}[$\star\star$]\label{exo:g7:areas:10}
一个半径为 $6$ cm 的圆柱形花瓶中装有水，水面高度为
$15$ cm。将全部水倒入一个空的长方体容器，
其底面为 $18$ cm $\times$ $12$ cm。
求容器中水面的高度。（体积相同，底面不同。）
\end{exercise}

\begin{exercise}[$\star\star$]\label{exo:g7:areas:11}
直径为 $30$ cm 的比萨售价 $9$ 欧元；直径为 $40$ cm 的
比萨售价 $14$ 欧元。求两张比萨的面积，以及每 $100$
cm$^2$ 的价格。哪一张更划算？
\end{exercise}

\begin{exercise}[$\star\star\star$]\label{exo:g7:areas:12}
直角三角形的两条直角边分别为 $6$ cm 和 $8$ cm，其
斜边长为 $10$ cm。先用两条直角边求面积，
再利用该面积，求相对于\emph{斜边}的高。
\end{exercise}

\section{问题：消失的正方形}

\begin{problem}[{周末问题 --- 一个著名的“证明”声称
$64 = 65$，借面积与斜率揭穿，
并以比萨与圆环作甜点}]
\label{pb:g7:areas:1}
一位魔术师把 $8 \times 8$ 的正方形剪成四块，
再将它们滑动、重新拼成一个 $5 \times 13$ 的长方形。
观众算道：$8 \times 8 = 64$，而
$5 \times 13 = 65$ --- 一个单位正方形竟
凭空出现了！本章的面积公式恰好是
能揭穿这一戏法的工具。本题接着以诚实的面积推理，
研究梯形、倾斜的堆叠和比萨的经济问题 ---
最后又回到魔术师，看他第二场
更为离奇的表演。

\medskip
\noindent\textbf{第一部分 --- 戏法的表演与揭穿。}
这个 $8 \times 8$ 正方形的四块是：两个直角
三角形，其直角边为 $8$ 和 $3$；两个直角梯形，
其平行边为 $5$ 和 $3$，高为 $5$。
\begin{enumerate}
  \item 求原正方形的面积，以及所声称的
        长方形的面积。用一句话说明其中的矛盾。
  \item 求四块图形各自的面积
        （见\cref{thm:g7:areas:triangle}；对梯形，可切开
        或利用\cref{exo:g7:areas:9}）。
  \item 将四块面积相加。在两个图形 --- 正方形与
        长方形 --- 中，这四块究竟能拼满哪一个？
  \item 在长方形的拼法中，三角形的斜边
        （上升 $3$、水平 $8$）应当与梯形的
        腰（上升 $2$、水平 $5$）连成一条笔直的“对角线”。
        检验对齐：上升 $3$、水平 $8$ 的斜率与
        上升 $2$、水平 $5$ 的是否相同？
        （见\cref{thm:g7:prop:cross}）？
  \item 解释多出的第 65 个单位正方形藏在哪里：沿假“对角线”的空隙是什么形状，
        它的精确面积是多少？
        为什么眼睛会忽略它？（这块面积拉成细条，
        铺在长方形的对角线上，平均宽度大约是多少？
        该对角线约长 $13$ 个单位。）
\end{enumerate}

\medskip
\noindent\textbf{第二部分 --- 诚实的面积推理。}
\begin{enumerate}[resume]
  \item 画两条相距 $4$ cm 的平行线，在其中一条直线上取长为 $6$ cm 的底边，
        并在该底边上作三个不同的三角形，顶点取在另一条直线
        上的不同位置。求这三个三角形的面积。
        为什么无论顶点怎样沿着直线滑动，
        这三个面积都保持相等？
  \item 一般地证明梯形面积公式：梯形的平行边分别为
        $B$ 与 $b$，高为 $h$，沿一条对角线切开，
        得到两个高均为 $h$ 的三角形。
        由此得出：
        \[
        A = \frac{(B + b) \times h}{2} .
        \]
  \item 将该公式应用于平行边为 $B = 9$ cm、$b = 5$ cm，
        高为 $h = 4$ cm 的梯形，求其面积。
  \item 木匠实际上也在用同一公式，只是换了说法：面积 $=$
        （两平行边的平均值）$\times$ 高。
        请说明这为何与上式是同一法则，并用这一方法重算第
        8 题。
  \item 将一摞码放整齐的扑克牌推成倾斜的堆叠。
        解释为何体积没有改变，这对棱柱体积公式中的高
        有何启示：
        $V = B \times h$（\cref{thm:g7:areas:volume}）中，
        对于斜的堆叠，应取哪一个高 --- 斜长还是
        竖直的那个？（对比
        \cref{exo:g7:areas:2}，这里升高一维。）
\end{enumerate}

\medskip
\noindent\textbf{第三部分 --- 圆环、比萨与魔术师的
回归。}
\begin{enumerate}[resume]
  \item 一个半径为 $3$ m 的圆形池塘，位于
        半径为 $5$ m 的圆形草坪正中央。求
        草地圆环的面积（精确值用 $\pi$ 表示，然后
        四舍五入到 m$^2$）。
  \item 求一个圆盘的半径，使其面积恰好等于
        该圆环的面积。（你现在得到的三个半径
        构成一组著名的勾股数，其故事在
        \cref{ch:g8:pythagoras}中讲述。）
  \item 哪一个比萨分量更多：直径为 $40$ cm 的一张比萨，还是
        直径为 $28$ cm 的两张比萨？请计算并比较。
  \item 要使比萨的面积\emph{翻倍}，其半径大约需增大为原来的
        多少倍？检验候选值 $1.4$（回忆
        面积随比例因子的\emph{平方}而变化，
        见\cref{pb:g6:propdata:1}），并说明精确的倍数 ---
        即平方等于 $2$ 的那个数 --- 在何处
        正式登场
        （见\cref{ex:g8:pythagoras:diagonal}）。
  \item 魔术师再次登场：他把一个 $13 \times 13$ 的正方形
        剪开，让数 $5$、$8$、$13$ 充当
        先前 $3$、$5$、$8$ 所扮演的角色，拼成一个
        $8 \times 21$ 的长方形。求两个面积：这一次
        消失的是什么？数 $3, 5, 8, 13, 21$ 恰是
        数列的连续项，该数列见
        \cref{pb:g7:fractions:1} --- 请写出魔术师的
        得失规律，以及其中的道理：为何面积
        从不说谎？
\end{enumerate}
\end{problem}
